\documentclass[11pt]{article}
\usepackage{amsmath,amsthm,amssymb}
\usepackage[margin=1in]{geometry}
\usepackage[hidelinks]{hyperref}
\newtheorem{theorem}{Theorem}
\newtheorem{corollary}[theorem]{Corollary}
\newtheorem{proposition}[theorem]{Proposition}
\newcommand{\Z}{\mathbb Z}
\newcommand{\res}{\mathbin{+_{R}}}
\title{Counterexamples to a conjecture on one-sided restricted sumsets}
\author{}
\date{September 8, 2026}
\begin{document}
\maketitle
\begin{abstract}
For finite integer sets $A,B$ and a forbidden relation $R\subseteq A\times B$, let $A\res B$ be the set of sums represented by pairs outside $R$. Ouyang conjectured that, when $|B|\leq |A|$ and every vertex of $B$ has forbidden degree at most $\Delta$, one has $|A\res B|\geq |A|+|B|-1-\lfloor 5\Delta/2\rfloor$. We give explicit counterexamples for every $\Delta=3k$, $k\geq1$. For every $N\geq14k$, our sets have cardinality $N$ and restricted sumset cardinality $2N-8k-1$, falling below the conjectured bound by $\lceil k/2\rceil$. Consequently a universal linear loss coefficient must be at least $8/3$. For $\Delta=3$ and equal summand sizes $N\geq14$, the construction attains the known lower bound $2N-9$, and hence determines the exact minimum.
\end{abstract}

\section{Introduction}
Let $A,B\subseteq\Z$ be finite, nonempty sets and let $R\subseteq A\times B$. The restricted sumset associated with the forbidden relation $R$ is
\[
 A\res B=\{a+b:a\in A,\ b\in B,\ (a,b)\notin R\}.
\]
For $b\in B$, put $d_R(b)=|\{a\in A:(a,b)\in R\}|$. We consider the one-sided condition $d_R(b)\leq\Delta$ for all $b\in B$. Degrees on $A$ are unrestricted.

Ouyang proved that, if $|B|\leq|A|$ and the one-sided degree condition holds, then
\begin{equation}\label{eq:known}
 |A\res B|\geq |A|+|B|-3\Delta
\end{equation}
\cite[Theorem 1.5(i)]{Ouyang}. Motivated by a construction in \cite[Theorem 1.9]{Ouyang}, Ouyang proposed the stronger inequality
\begin{equation}\label{eq:conjecture}
 |A\res B|\geq |A|+|B|-1-\left\lfloor\frac{5\Delta}{2}\right\rfloor
\end{equation}
\cite[Conjecture 1.10]{Ouyang}. All numbered references to that paper refer to its September 4, 2025 author version, arXiv:2503.09121v3.

We disprove \eqref{eq:conjecture} by a family whose forbidden degree is unbounded. The construction uses an interval with three holes as $A$, an interval as $B$, and four intervals of deleted sums. The holes near the endpoints of $A$ control the degree when both interior intervals of deleted sums meet the same translate of $A$. Exact intersection counts show that the resulting degree is at most $3k$, while the loss from $|A|+|B|-1$ is $8k$.

Throughout, $[x,y]$ denotes the integer interval $\{x,x+1,\ldots,y\}$ for integers $x\leq y$.

\begin{theorem}\label{thm:construction}
Let $k,N,U$ be integers satisfying
\begin{equation}\label{eq:parameters}
 k\geq1,\qquad N\geq14k,\qquad 8k\leq U\leq N-6k.
\end{equation}
Define
\begin{align*}
 A={}&[0,N+4k-1]\setminus\bigl([k,2k-1]\cup[U,U+2k-1]\nobreak\\[-2pt]
 &\hspace{43mm}\cup[N+2k,N+3k-1]\bigr),\\
 B={}&[0,N-1],
\end{align*}
and put
\begin{align*}
 D_0&=[0,4k-1],& D_1&=[U,U+2k-1],\\
 D_2&=[U+N-1,U+N+2k-2],&
 D_3&=[2N-1,2N+4k-2].
\end{align*}
Let $D=D_0\cup D_1\cup D_2\cup D_3$ and
\[
 R=\{(a,b)\in A\times B:a+b\in D\}.
\]
Then
\begin{gather*}
 |A|=|B|=N,\qquad \max_{b\in B}d_R(b)=3k,\\
 A\res B=[0,2N+4k-2]\setminus D,\qquad
 |A\res B|=2N-8k-1.
\end{gather*}
\end{theorem}

\begin{corollary}\label{cor:failure}
Conjecture 1.10 of \cite{Ouyang} is false. For each $k\geq1$, Theorem~\ref{thm:construction} gives examples with $\Delta=3k$ whose restricted sumset has $\lceil k/2\rceil$ fewer elements than the proposed lower bound.
\end{corollary}
\begin{proof}
The proposed lower bound is $2N-1-\lfloor15k/2\rfloor$, whereas the actual cardinality is $2N-8k-1$. Their difference is
\[
 8k-\left\lfloor\frac{15k}{2}\right\rfloor
 =\left\lceil\frac{k}{2}\right\rceil>0.
\]
\end{proof}

\begin{corollary}\label{cor:coefficient}
Suppose real constants $c,K$, independent of $A,B,\Delta$, satisfy
\[
 |A\res B|\geq |A|+|B|-1-c\Delta-K
\]
for all finite, nonempty $A,B\subseteq\Z$ with $|B|\leq|A|$ and all relations of degree at most $\Delta$ on $B$. Then $c\geq8/3$.
\end{corollary}
\begin{proof}
Apply Theorem~\ref{thm:construction} with $N=14k$ and $U=8k$. The claimed universal bound implies $8k\leq3ck+K$. Divide by $k$ and let $k$ tend to infinity.
\end{proof}

\section{Proof of the construction}
\begin{proof}[Proof of Theorem~\ref{thm:construction}]
The three holes in $A$ are disjoint and have total cardinality $4k$. Hence $|A|=N$, and clearly $|B|=N$. Both endpoints $0$ and $N+4k-1$ belong to $A$. The largest difference between consecutive elements of $A$ is $2k+1\leq N$. Thus the intervals $a+B=[a,a+N-1]$, as $a$ runs through $A$ in increasing order, overlap or are adjacent. It follows that
\begin{equation}\label{eq:full}
 A+B=[0,2N+4k-2].
\end{equation}
The parameter inequalities imply that $D_0,D_1,D_2,D_3$ are disjoint, lie in this interval, and have lengths $4k,2k,2k,4k$, respectively. Consequently $|D|=12k$.

It remains to prove the degree assertion. For $b\in B$, define
\[
 d_i(b)=|(A+b)\cap D_i|\quad(0\leq i\leq3),\qquad
 d(b)=\sum_{i=0}^3d_i(b)=d_R(b).
\]
The interval $D_0$ can contribute only when $0\leq b\leq4k-1$, and $D_3$ can contribute only when $N-4k\leq b\leq N-1$. The two interior intervals can both contribute only when
\[
 U-4k\leq b\leq U+2k-1.
\]
Indeed, meeting $D_2$ requires $N+4k-1+b\geq U+N-1$, while meeting $D_1$ requires $b\leq U+2k-1$. The three ranges
\begin{equation}\label{eq:ranges}
 [0,4k-1],\qquad [U-4k,U+2k-1],\qquad [N-4k,N-1]
\end{equation}
are pairwise disjoint: $U-4k\geq4k$ and $U+2k-1\leq N-4k-1$. Outside their union, neither endpoint interval contributes and at most one interior interval contributes. Its length is $2k$, so $d(b)\leq2k$ there.

\medskip\noindent\textit{The left endpoint range.}
Suppose $0\leq b\leq4k-1$. Only $D_0$ and $D_1$ contribute. The first count is the number of elements in $[0,4k-1-b]$ after removing $[k,2k-1]$. The second is the number of elements of $A$ in $[U-b,U+2k-1-b]$. This interval can meet the central hole $[U,U+2k-1]$ but no other hole. Direct counting gives
\begin{equation}\label{eq:lefttable}
\begin{array}{c|c|c|c}
 \text{range of }b&d_0(b)&d_1(b)&d(b)\\\hline
 0\leq b\leq2k-1&3k-b&b&3k\\
 2k\leq b\leq3k-1&k&2k&3k\\
 3k\leq b\leq4k-1&4k-b&2k&6k-b
\end{array}
\end{equation}
Every entry in the last column is at most $3k$.

\medskip\noindent\textit{The middle range.}
Write $b=U-4k+t$ with $0\leq t\leq6k-1$. Only $D_1$ and $D_2$ contribute. To compute $d_1(b)$, intersect $A$ with
\[
 D_1-b=[4k-t,6k-1-t].
\]
This interval meets only the left endpoint of $A$ and the hole $[k,2k-1]$; its upper endpoint is less than $U$. For the second count,
\[
 D_2-b=[N+4k-1-t,N+6k-2-t].
\]
Set $L=N+4k-1$. Reflection $x\mapsto L-x$ takes this interval to $[t-2k+1,t]$. The reflected set $L-A$ has the endpoint hole $[k,2k-1]$ and its central hole begins at $N+2k-U\geq8k$. Since $t\leq6k-1$, no other hole contributes to this count. The exact intersection sizes are therefore
\begin{equation}\label{eq:middletable}
\begin{array}{c|c|c|c}
 \text{range of }t&d_1(b)&d_2(b)&d(b)\\\hline
 0\leq t\leq k-1&2k&t+1&2k+t+1\\
 k\leq t\leq2k-1&2k&k&3k\\
 2k\leq t\leq3k-1&4k-t&k&5k-t\\
 3k\leq t\leq4k-1&k&t-2k+1&t-k+1\\
 4k\leq t\leq5k-1&k&2k&3k\\
 5k\leq t\leq6k-1&6k-t&2k&8k-t
\end{array}
\end{equation}
Again every entry in the last column is at most $3k$.

\medskip\noindent\textit{The right endpoint range.}
The left endpoint calculation also handles this range by reflection. To make the parameter change explicit, write $A_U,D_U$ for the constructed sets. With
\[
 U'=N+2k-U,\qquad M=2N+4k-2,
\]
we have
\[
 L-A_U=A_{U'},\qquad M-D_U=D_{U'},\qquad (N-1)-B=B.
\]
Moreover, $8k\leq U'\leq N-6k$. The map
\[
 (a,b)\longmapsto(L-a,N-1-b)
\]
therefore carries the forbidden relation for $U$ to the forbidden relation for $U'$. For $b\in[N-4k,N-1]$, the new coordinate $N-1-b$ lies in $[0,4k-1]$, where \eqref{eq:lefttable} applies. This proves $d(b)\leq3k$ in the right endpoint range.

Together these calculations establish $d_R(b)\leq3k$ for every $b\in B$. Equality holds at $b=0$ by \eqref{eq:lefttable}. Finally, by the definition of $R$, every representation of a sum in $D$ is forbidden and every representation of a sum outside $D$ is permitted. Using \eqref{eq:full}, we obtain
\[
 A\res B=(A+B)\setminus D=[0,2N+4k-2]\setminus D.
\]
Its cardinality is $(2N+4k-1)-12k=2N-8k-1$, as required.
\end{proof}

\section{Exact sharpness at degree three}
For completeness, we recall the elementary lower-bound argument underlying \eqref{eq:known}. This is the argument of \cite[Theorem 1.5(i)]{Ouyang}.

\begin{proposition}[Ouyang]\label{prop:known}
Let $A,B\subseteq\Z$ be finite, nonempty sets, with $|B|\leq|A|$, and let $\Delta\geq1$ be an integer. If $R\subseteq A\times B$ has degree at most $\Delta$ on $B$, then $|A\res B|\geq |A|+|B|-3\Delta$.
\end{proposition}
\begin{proof}
Write $A=\{a_1<\cdots<a_m\}$ and $B=\{b_1<\cdots<b_n\}$. If $a_1$ has forbidden degree at most $\Delta$, the strictly increasing chain of sums along
\[
 (a_1,b_1),\ldots,(a_1,b_n),(a_2,b_n),\ldots,(a_m,b_n)
\]
has $m+n-1$ elements and at most $2\Delta$ forbidden pairs. It supplies at least $m+n-1-2\Delta\geq m+n-3\Delta$ permitted sums.

Otherwise, $a_1$ has degree at least $\Delta+1$. Since
\[
 \sum_{a\in A}d_R(a)=\sum_{b\in B}d_R(b)\leq n\Delta\leq m\Delta,
\]
there is some $a_j$ with $j\geq2$ whose degree is at most $\Delta-1$. Consider the chain along the first column up to $a_j$, then along the row of $a_j$, then along the last column:
\[
 (a_1,b_1),\ldots,(a_j,b_1),
 (a_j,b_2),\ldots,(a_j,b_n),
 (a_{j+1},b_n),\ldots,(a_m,b_n).
\]
It contains $m+n-1$ strictly increasing sums. The number of forbidden pairs is at most $\Delta+(\Delta-1)+\Delta=3\Delta-1$, since the two columns have degree at most $\Delta$. Thus at least $m+n-3\Delta$ sums remain.
\end{proof}

\begin{corollary}\label{cor:exact}
For every integer $N\geq14$,
\[
 \min |A\res B|=2N-9,
\]
where the minimum is over finite integer sets $A,B$ of cardinality $N$ and relations $R\subseteq A\times B$ of degree at most three on $B$.
\end{corollary}
\begin{proof}
Proposition~\ref{prop:known} gives the lower bound. Theorem~\ref{thm:construction}, with $k=1$ and any $8\leq U\leq N-6$, attains it.
\end{proof}

For example, take $k=1$, $N=20$, and $U=14$. Then
\begin{align*}
 A&=\{0\}\cup[2,13]\cup[16,21]\cup\{23\},\\
 B&=[0,19],\\
 D&=\{0,1,2,3,14,15,33,34,39,40,41,42\}.
\end{align*}
The degree sequence on $B$, in increasing order, is
\[
 (3,3,3,3,2,2,2,2,2,2,3,3,3,3,3,3,3,3,3,3).
\]
The restricted sumset has $31$ elements, whereas \eqref{eq:conjecture} requires $32$.

The construction leaves a gap in the possible universal leading coefficient: Corollary~\ref{cor:coefficient} requires at least $8/3$, while \eqref{eq:known} supplies $3$. We do not determine that coefficient, classify the equality cases at degree three, or assert a corresponding counterexample under a degree restriction on both sides.

\begin{thebibliography}{9}
\bibitem{Ouyang}
Minghui Ouyang,
\emph{On restricted sumsets with bounded degree relations},
Mathematika \textbf{71} (2025), no.~4, e70045.
\href{https://doi.org/10.1112/mtk.70045}{doi:10.1112/mtk.70045}.
Author version: \href{https://arxiv.org/abs/2503.09121v3}{arXiv:2503.09121v3}, September 4, 2025.
\end{thebibliography}
\end{document}
